% Generated draft: requires mathematical and source review.
\documentclass[11pt]{article}
\usepackage[margin=25mm]{geometry}
\usepackage{fontspec}
\setmainfont{Latin Modern Roman}
\usepackage{amsmath,amssymb,mathtools,unicode-math}
\setmathfont{Latin Modern Math}
\usepackage{enumitem,needspace,xurl,hyperref}
\setlength{\parindent}{0pt}
\setlength{\parskip}{5pt}
\setlength{\emergencystretch}{3em}
\newcommand{\N}{\mathbb{N}}
\newcommand{\Z}{\mathbb{Z}}
\newcommand{\Q}{\mathbb{Q}}
\newcommand{\R}{\mathbb{R}}
\newcommand{\C}{\mathbb{C}}
\newcommand{\F}{\mathbb{F}}
\newcommand{\A}{\mathbb{A}}
\newcommand{\PP}{\mathbb{P}}
\newcommand{\E}{\mathbb{E}}
\newcommand{\eps}{\varepsilon}
\newcommand{\dd}{\,\mathrm{d}}
\begin{document}
\section*{UnsolvedMath ID: 800578}
\textbf{Source problem code:} catalogue entry
\subsection{Definitions and mathematical statement}
Give each set \(X\subseteq\R\) its subspace topology. An open \(\omega\)-cover of \(X\) is a family \(\mathcal U\) of open subsets of \(X\) such that \(X\notin\mathcal U\) and, for every finite \(F\subseteq X\), some \(U\in\mathcal U\) satisfies \(F\subseteq U\). Write \(\Omega(X)\) for the family of countable open \(\omega\)-covers of \(X\). This uses the standard countable-open-cover convention; the supplied record does not explicitly define \(\Omega\).

The property \(\mathsf S_{\mathrm{fin}}(\Omega,\Omega)\) means that, for every sequence \((\mathcal U_n)_{n\in\N}\) with \(\mathcal U_n\in\Omega(X)\), there exist finite subfamilies \(\mathcal F_n\subseteq\mathcal U_n\) such that
\[
\bigcup_{n\in\N}\mathcal F_n\in\Omega(X).
\]
Here the union is the family of all selected cover members, rather than the family of their finite unions within individual covers.

Let \(\N=\{0,1,2,\ldots\}\). For \(f,g\in\N^{\N}\), write \(f\leq^{*}g\) if \(f(n)\leq g(n)\) for all sufficiently large \(n\). The dominating number is
\[
\mathfrak d=\min\left\{|D|:D\subseteq\N^{\N}\ \text{and}\ (\forall f\in\N^{\N})(\exists g\in D)\ f\leq^{*}g\right\}.
\]

The supplied existence question asks whether ZFC proves
\[
\exists X\subseteq\R\quad\bigl(|X|=\mathfrak d\ \text{and}\ X\text{ satisfies }\mathsf S_{\mathrm{fin}}(\Omega,\Omega)\bigr).
\]
The phrase ``in ZFC'' is interpreted here as asking for a ZFC theorem of existence.

The record also contains an inserted update reporting a negative answer, with Laver's model as a counterexample and reference label [BBC]. It further reports an affirmative answer if choosing one element from each cover is replaced by choosing a union of two elements from each cover, with reference label [MHP]. These reports are preserved as supplied qualifications, not as independently verified conclusions. The one-element/two-element language does not specify a modification of the displayed finite-selection property, so the update's relationship to this existence question remains unresolved. The registry separately labels the record ``open''; no resolution of that status conflict is asserted.
\subsection{Short English statement}
Does ZFC prove that there is a set of reals of cardinality \(\mathfrak d\) satisfying \(\mathsf S_{\mathrm{fin}}(\Omega,\Omega)\)? This property requires finite selections from every sequence of open \(\omega\)-covers whose combined family is again an \(\omega\)-cover. The supplied negative-answer update and its one-element/two-element qualification have an unclear relationship to this question and require source review.
\subsection{Sources}
\begin{itemize}
\item Local UnsolvedMath registry, permanent ID \(800578\), source problem code \texttt{catalogue entry}, AMR Open Problem Lists. The supplied record lists Tsaban as proposer, gives the year 2006, and labels the problem ``open''.
\item Tsaban, \emph{Selection Principles and Special Sets of Reals: Open Problems} (2006), Problem 4. Supplied source URL: \url{https://arxiv.org/abs/math/0609601}. The record reports access on 2026-10-16 and extraction by \texttt{arxiv\_latex}; these are supplied provenance details, not a claim of independent access or verification.
\item The title and background supply the attribution/context label [SFH], without its expansion or bibliographic details.
\item The inserted update cites [BBC] for its reported negative answer and Laver-model counterexample, and [MHP] for its reported affirmative answer involving unions of two cover elements. No bibliographic entries or links for these labels are supplied. Their identities and applicability to the displayed question remain unverified.
\end{itemize}
\end{document}
